\documentclass[11pt]{article}
\usepackage[margin=0.85in]{geometry}
\usepackage[T1]{fontenc}
\usepackage{amsmath,amssymb,booktabs,longtable,microtype,xurl,hyperref}
\setlength{\emergencystretch}{3em}
\hypersetup{colorlinks=true,urlcolor=blue,linkcolor=blue}
\title{Synchronizing the Manuscript with the Current Estimator and Controller}
\author{Pan Dynamics Collision Avoidance Research}
\date{October 2, 2026}
\begin{document}
\maketitle

\section{Scope}
This is a manuscript revision. No controller, estimator, configuration or
test source is changed. \path{paper/manuscript.tex} had last been
synchronized on September 19, 2026 (commit \path{e64d5dd}, with a
single-target edit in \path{c152323}) and described the retired
joint-support controller: one convex program with jointly optimized
separating angles, a modal terminal set, an exit deadline and a conditional
recursive-feasibility proposition on a held affine plant. Since then the
controller was replaced and the estimator gain synthesis changed.

The revision describes the algorithms committed at
\path{5692e1551241b568bcfb300093814373c674a412}. The controller sources of
that commit are those of \path{a29ccb6ff2856fe3df33efff4b1cb8619e523658},
which produced the closed-loop records used here; the hashes of
\path{controller/solvePredictiveControl.m} and
\path{controller/predictiveSafetyGeometry.m} at the reviewed commit equal
the hashes stored in
\path{report/SHARMA_TIMED_FLOW_20261002/provenance.json}.

When the revision was written, the working tree held uncommitted edits of a
concurrent session to the controller (within-hold relinearization with a
nonlinear model-agreement test, and a closed-form distance-dual
construction), its configuration, tests and documents. Those edits have no
recorded result. The manuscript does not describe them as implemented; it
names both as next steps. All controller facts were read from an export of
the reviewed commit (\texttt{git archive HEAD controller config scripts}),
not from the working tree. \path{provenance.json} lists the fifteen
concurrently modified paths that this change leaves untouched.

\section{What changed in the manuscript}
The section order is kept. The literature discussion of the introduction
and the perception section are byte-identical to the preceding revision, and
the audit checks both. The estimator derivations before the target-tracker
gain statement receive small clarifications only (43 changed lines of 324):
the definition of the ego-frame acceleration, the link between the
relative-position measurement and the perception output, the description of
the reference observer's switched yaw stage, the wording of the
single-track assumption and of one lemma hypothesis, the target speed
bounds, the symbols of the Lipschitz proof, the component constants of the
ISS theorem, the yaw-error comparison and the first-use definition of RK4.
The word ``radar'' is replaced by ``relative-position measurement''
because the perception section describes a camera--LiDAR pipeline.

\begin{itemize}
\item \textbf{Title, abstract, keywords, last introduction paragraph.} The
second half of the title changes from ``Covariant Ego--Target Estimation and
Conditional Control Analysis'' to ``Covariant Ego--Target Estimation and
Two-Stage Predictive Control''. The abstract (250 words) and the closing
introduction paragraph now state the certified estimator, the two-stage
controller, the measured outcome (twelve of fourteen runs collision-free
with recovery, one collision, one numerical interruption, call time above
the hold), the absence of a controller certificate, and that the two
components are not evaluated together.
\item \textbf{System overview.} The longitudinal bias term and the
regularized tire denominator are removed, in agreement with
\path{nonlinearBicycleModel.derivative}. The text states that the reference
is straight or of constant curvature, that the controller is
certainty-equivalent, and that it imposes no road-boundary constraint. The
architecture figure is redrawn (\path{system_architecture_v11}).
\item \textbf{Estimator.} The bounded-real shape program is replaced by the
noise-variance program of \path{synthesizeTargetTrackerCertificate.m}
(dissipation inequality, observability-Gramian inequality and variance
epigraph in one metric), followed by the unchanged free-metric bandwidth
search. The enclosure subsection states that the controller does not yet
consume the enclosure.
\item \textbf{Control design.} Rewritten: RK4 hold map and its tangents,
closed-form obstacle flow, anchor and trust boxes, shifted-rollout and
moving-flow initialization, distance-dual rows with fixed multipliers,
free-pose terminal core, the path-guidance law, the nominal cost-to-go
function and its local convex model, the two cone programs, restoration and
execution rules, and an explicit scope of guarantees. Three formal
statements are given with proofs: the fixed-multiplier lower bound with its
closed-form multipliers and its degenerate overlap case, the Schur
reduction of the freely placed core, and the telescoping decrease of the
cost-to-go under its construction feedback, including the positive
invariance of its domain.
\item \textbf{Results.} Current estimator design values; the held-out
comparison with the reproduced, per-scenario tuned reference observer,
including the per-scenario synthesis of the estimator gains and the metrics
that favor the reference; the fourteen-scenario closed-loop campaign with a
parameter table, an outcome table and a figure; and the failure mechanisms
of the collision and of the interruption.
\item \textbf{Discussion and conclusion.} Rewritten to separate the
continuous-time estimation certificate from the uncertified controller and
to list the favorable premises of the closed-loop trials.
\item \textbf{Declarations.} The data-availability statement names the
external raw traces and diagnostic helpers, and the AI-assistance
disclosure names the tool and the parts it drafted.
\end{itemize}

Six references are added: Diehl, Bock and Schl\"oder (real-time iteration),
Jadbabaie, Yu and Hauser, Rawlings, Mayne and Diehl (terminal decrease),
Boyd et al.\ (Gramian inequality), L\"ofberg (YALMIP) and Sturm (SeDuMi).
Five were checked against Crossref metadata by DOI on October 2, 2026, and
the book of Rawlings, Mayne and Diehl against its publisher record
(Nob Hill Publishing, second edition, 2017, ISBN 978-0-9759377-3-0). The
Crossref container title of the YALMIP record is wrong; the entry uses the
CACSD 2004 proceedings named by the DOI. Two preprint entries are replaced
by their published versions, both checked against Crossref on the same day:
Zhang, Liniger and Borrelli, \emph{IEEE Transactions on Control Systems
Technology} 29(3):972--983, 2021 (DOI \path{10.1109/TCST.2019.2949540}),
and Zeng, Zhang, Li and Sreenath, 2021 American Control Conference,
pp.~3856--3863 (DOI \path{10.23919/ACC50511.2021.9482626}). The bibliography
file holds 44 entries, of which the manuscript cites 35.

\section{Evidence used and work actually run}
\paragraph{Records read, not rerun.}
\path{report/SHARMA_TIMED_FLOW_20261002.tex} and its
\path{comparison.json}; \path{report/CONTROLLER_FAILURE_CAUSES_20261002.tex}
and its \path{findings.json};
\path{report/TARGET_NOISE_SHAPE_GAINS_20261001.tex} with its design,
benchmark and audit tables; \path{report/SHARMA_ESTIMATOR_AUDIT_20260930.tex};
the free-pose terminal, encounter-range and unified-CLF reports; and the
controller documents \path{PCBF_CLF_ARCHITECTURE.md} and
\path{NOMINAL_CLF.md} at the reviewed commit.

\paragraph{Computations run for this revision.}
\begin{enumerate}
\item \texttt{evaluateEstimatorDesign}: the committed
\path{synthesizeNrmmObserverGains} on \path{nrmmTrackingConfig}, MATLAB
R2026a Update 3, one computational thread, YALMIP and SeDuMi from
\path{solver/}. Output \path{estimator-design.json}. It reproduces the
default-domain row of the recorded design table (bandwidth 2.1042, gains
7.2977, 26.2540, 40.9423, radii 1.7144 m, 7.554 m/s, 16.52 m/s$^2$).
\item \texttt{evaluateControllerConstants}: terminal-core construction,
cruise trims and the spectral radius of the nominal feedback at 8 and
15 m/s, evaluated on the export of the reviewed commit. Output
\path{controller-constants.json}. The certified core radii are
$2.4\times10^{-4}$ and $9.8\times10^{-4}$ with contraction bounds 0.9996
and 0.9987.
\item \path{deriveClosedLoopStatistics.py}: statistics of the recorded
campaign. It verifies the SHA-256 of each of the fourteen raw traces against
\path{comparison.json}, recomputes the body distance with the independent
Python geometry, and writes \path{closed-loop-statistics.json}. No
simulation is run. The raw traces remain in
\path{/home/zai/.cache/collisionAvoidance/sharma-flow-20261002/current-campaign/}.
\item \path{figure-sources/closed_loop_results_v01.py}: the results figure,
drawn from the same hash-verified traces.
\item \path{checkDualLemma.py}: a numerical test of the three parts of the
fixed-multiplier lemma and of the first-order vertex-rotation remainder on
40,000 random rectangle pairs (32,756 disjoint, 7,244 overlapping) with
geometry written independently of the controller. It passes. It tests the
statement and is not a proof.
\item \path{auditManuscript.py}: recomputes every quoted number from its
record and requires the formatted value in the manuscript; checks labels,
references, citation keys, environment balance, the absence of
retired-controller vocabulary and the spans that the revision leaves
unchanged.
\item \path{writeProvenance.py}: source hashes at the reviewed commit,
record hashes, tool versions, commands and output hashes, written to
\path{provenance.json}.
\end{enumerate}

No closed-loop simulation, estimator trial or unit-test suite was run. The
closed-loop statistics that do not appear in the earlier reports (call-time
distribution, lateral excursions, corridor exits, counts of positive-slack
samples, control-Lyapunov decrease statistics) are derived from the recorded
traces by item 3.

\section{Derived statistics new to this revision}
\begin{itemize}
\item 3,708 controller calls; all exceed the 50 ms hold; median 0.107 s,
95th percentile 0.360 s, maximum 0.976 s (first call of a resumed process;
0.902 s otherwise); 2,252 exceed 0.1 s; solver share of call time 0.67.
\item In seven of the twelve successful runs the smallest sampled distance
is below the 0.10 m node margin (minimum 0.033 m). Four of these seven
runs executed at least one plan with a positive primary optimum.
\item Largest lateral deviation 2.5 to 9.9 m; in twelve of fourteen runs
the body leaves a corridor of half-width 4 m about the path. Largest
steering angle 0.59 rad, largest yaw rate 0.91 rad/s, force ratio between
$-0.68$ and $0.56$.
\item A positive primary optimum in 22 samples of five runs, between 0.80
and 1.50 s. A fresh anchor is formulated at 27 samples and restores a zero
optimum at 7; the 22 samples are the 20 unsuccessful ones and two at which
the optimum equals its lower bound. 113 executed solutions carry a
non-convergence exit flag.
\item Of 3,694 consecutive measured-state pairs, 2,636 have zero modeled
control-Lyapunov slack; in 154 of them the requested decrease is not
achieved and in 36 the value increases (at most 0.034).
\end{itemize}

\section{Review passes and corrections}
Three independent read-only review passes were run by separate AI reviewer
agents. Each received the manuscript and the records, not the drafting
notes, and none edited a file. Their findings were checked against the
sources before any change was made.

\paragraph{Controller facts against the reviewed commit.} Corrected: the
restoration stage of the collision (it fails once both the shifted and the
fresh anchor overlap; the fresh-anchor counts and the 0.90 s sample are now
stated); the lexicographic tolerance in the stated separation bound; the
order of the two terminal separation constructions; the feedback that
completes a fresh anchor; the trim used to complete a shifted plan; the
description of the turning-crossing scenario; the solver and integration
tolerances; the statement that the passing side is fixed by the
initialization; and the mismatch between the closed-loop scenarios and the
estimator's certified domain. One reviewer count was not adopted: the
records show four, not three, sub-margin successes with executed
relaxations, and the manuscript says four.

\paragraph{Estimator facts against the records.} Corrected: the estimator
gains of the comparison are synthesized anew for each scenario's declared
domain, which the draft had not disclosed; a rounding slip (0.054 to
0.053 m/s); the agreement of the covariant and companion chains under
identical gains (three of four constant-motion scenarios, not all); the
description of the reference's yaw stage; ``the shape of the gains is
obtained from'' instead of ``the gains solve''; the sign convention of the
noise term; and ``4 to 5\%'' for the small position advantages. Added: the
asymmetries of the comparison (the reference does not use GNSS velocity,
its selected candidate lies on the grid boundary in five of six scenarios,
the stress cells have five seeds, and the noise-variance criterion was
adopted after an earlier comparison on the same test seeds was known) and
the less favorable metrics (acceleration lower in two of six scenarios,
whole-run position in none, and the noise-free results).

\paragraph{Internal consistency and prose.} Corrected: clashing symbols in
the control section; the statement of the fixed-multiplier lemma and the
strengthened cost-to-go proposition; an explicit guidance-law equation with
its parameters in the parameter table; descriptions and consistent names of
the six estimator scenarios; consistent names for the two gain shapes;
number formats; units of the primary optimum; acronym definitions in the
rewritten text; one displayed equation layout; and the two published
reference versions. The guidance-law equation was added after the first
pass and was afterwards compared line by line with
\path{nonlinearBicycleModel.nominalFeedback} at the reviewed commit. The
final proofread renamed three guidance-law symbols that clashed with
earlier notation.

\section{Review items left open}
These items were raised and deliberately not changed. They need the
author's decision.
\begin{itemize}
\item The first half of the title, ``Integrated Lane Keeping and Vehicle
Collision Avoidance'', is kept although the estimator and the controller
are evaluated separately. The reviewer suggested a title that states the
separate evaluation.
\item The estimator comparison table reports means with a mark where the
paired bootstrap interval contains zero; standard deviations or intervals
for every cell are not added.
\item The unchanged introduction and perception text keep acronyms without
definition (MPC, QP, YOLO, LiDAR) and one British spelling. They are left
because those spans are kept byte-identical.
\item The perception section is not used or evaluated by any reported
experiment, and the error enclosure of the estimator is not consumed by the
controller. The manuscript states both.
\item Nine bibliography entries are uncited, the \texttt{multirow} package
and one macro are unused, and one symbol of the vehicle-dynamics figure is
not used in the text.
\end{itemize}

\section{Limits of this revision}
The manuscript is a research draft. Author, affiliation, funding,
competing-interest, contribution and repository-release declarations are
still placeholders. The estimator and the controller are evaluated
separately and with different vehicle parameter sets; the manuscript says
so. The enclosure-containment and outage campaigns of earlier estimator
gains were not repeated and are no longer quoted. The figure assets under
\path{paper/figures/} are untracked, as before; the sources of the two new
assets are kept in this bundle. \path{CLAUDE.md} still describes the retired
controller and was not edited in this task.

\section{Build and checks}
\begin{itemize}
\item \texttt{latexmk -pdf -interaction=nonstopmode -halt-on-error
manuscript.tex} in \path{paper/} exits with status 0 and produces 19 pages.
The log contains no LaTeX or package warning, no undefined reference or
citation and no overfull box; nine underfull-box notices remain.
\item \path{auditManuscript.py} with the preceding revision
\path{5692e1551241b568bcfb300093814373c674a412}: 128 of 128 checks pass
(\path{audit.json}). In a negative test on a copy of the manuscript with
five altered numbers, the five corresponding checks fail and the exit
status is nonzero.
\item \path{checkDualLemma.py} with 40,000 samples passes.
\item \texttt{git diff --check} reports nothing for the committed files.
\end{itemize}

\section{Files}
\path{paper/manuscript.tex}, \path{paper/references.bib},
\path{paper/README.md}, this report and the bundle
\path{report/MANUSCRIPT_ALGORITHM_SYNC_20261002/}:
\path{auditManuscript.py}, \path{audit.json},
\path{deriveClosedLoopStatistics.py}, \path{closed-loop-statistics.json},
\path{evaluateEstimatorDesign.m}, \path{estimator-design.json},
\path{evaluateControllerConstants.m}, \path{controller-constants.json},
\path{checkDualLemma.py}, \path{writeProvenance.py}, \path{provenance.json}
and \path{figure-sources/}. Concurrent controller, configuration, test,
estimator-document and report-index edits, agent instruction files, figure
binaries, the compiled manuscript, literature PDFs and solver dependencies
are not part of this change.
\end{document}
